\section{研究局限性}

本文的实验对象固定在树莓派 5、SO-101 机械臂和 ACT 策略模型上，
因此结论首先服务于 ARM CPU 边缘机器人这一类受限平台。
不同机械臂的通信协议、相机驱动、策略模型结构和任务动作尺度可能改变各阶段占比，
尤其是在模型更小、舵机总线更慢或相机分辨率更高的配置中，
瓶颈可能从推理运行时转移到硬件 I/O 或图像预处理路径。
因此，本文得到的具体数值不应直接外推到所有机器人平台，
更适合作为一套跨层性能分析方法和基础软件优化思路的案例验证。

其次，本文已经对异步推理进行了实测验证，但对推理运行时本体的优化仍然有限。
因此，本文能够说明当前系统为什么慢、异步调度能恢复多少吞吐，
但还不能给出降低单次 ACT 推理耗时的最终工程方案。

最后，本文的动作质量评价仍以操作者现场观察为主，尚未建立独立于主观判断的任务质量指标。
异步推理中的 stall 比例、过期帧比例和有效 FPS 可以反映控制流水线的运行状态，
但它们并不能完全替代抓取成功率、末端轨迹误差、动作平滑度和安全裕量等机器人任务指标。
后续若要将本文方法推广为更通用的基础软件评测框架，
仍需把系统级性能指标与任务级行为质量进一步关联起来。
